% !TEX root = relatorio.tex
%
% Capítulo 3 — Análise e Requisitos
%
\chapter{Análise e Requisitos} \label{cap:requisitos}

A definição clara dos requisitos de um sistema é um dos fatores mais determinantes para o sucesso de um projeto de engenharia de software~\cite{sommerville2016}.
Este capítulo apresenta a análise dos utilizadores da plataforma Play4Change, os requisitos funcionais e não funcionais identificados, e os principais casos de uso que guiaram as decisões de desenho e implementação.

% -----------------------------------------------------------
\section{Identificação dos Utilizadores} \label{sec:utilizadores}
% -----------------------------------------------------------

A identificação dos utilizadores é o ponto de partida de qualquer processo de desenvolvimento centrado no utilizador (\textit{User-Centered Design})~\cite{norman1988}, e levou, no contexto do Play4Change, à identificação de dois perfis de utilizador com necessidades, objetivos e contextos de uso distintos.

\subsection{Administrador}

O \textbf{Administrador} é o perfil responsável pela criação e gestão de conteúdo pedagógico --- tipicamente um formador ou responsável de uma instituição parceira da rede U!REKA, sem necessidade de \textit{expertise} em IA ou engenharia de software. Interage com a plataforma exclusivamente através do \textbf{cliente web}: submete PDF ou URLs como fontes de conhecimento, revê os tópicos e desafios gerados pela IA, e gere os Learners registados. Esta separação --- produção de conteúdo versus consumo e progressão --- é consistente com os modelos de \textit{Intelligent Tutoring Systems} (ITS)~\cite{anderson1985}.

\subsection{Learner}

O \textbf{Learner} é o destinatário da experiência de aprendizagem. O seu perfil é deliberadamente amplo: qualquer cidadão que pretenda desenvolver as suas competências digitais no referencial DigComp~2.2, independentemente do nível de partida ou do contexto institucional em que se insere.

O Learner interage com a plataforma através do \textbf{cliente móvel}, numa experiência desenhada para ser discreta e integrada na rotina diária. O modelo de um \textbf{desafio por dia} não é arbitrário: sessões curtas e regulares são significativamente mais eficazes do que sessões longas e esporádicas (\textit{spaced repetition})~\cite{ebbinghaus1885,cepeda2006}, e a Teoria da Autodeterminação de Deci e Ryan mostra que a motivação intrínseca é sustentada por três necessidades psicológicas básicas --- \textit{autonomia}, \textit{competência} e \textit{relatedness} (ligação social a outras pessoas, não apenas a um propósito abstrato)~\cite{deci2000} --- fatores a que este modelo responde parcialmente: a autonomia e a competência são endereçadas diretamente pelo desenho do desafio diário, mas a \textit{relatedness} depende das funcionalidades sociais ainda por implementar (Secção~\ref{sec:trabalho-futuro}). A escolha do canal móvel segue a justificação já apresentada na Secção~\ref{sec:formal-informal}.

% -----------------------------------------------------------
\section{Requisitos Funcionais} \label{sec:rf}
% -----------------------------------------------------------

Os requisitos funcionais descrevem os comportamentos que o sistema deve exibir em resposta a determinadas entradas ou condições~\cite{sommerville2016}.
A Tabela~\ref{tab:rf} apresenta os requisitos funcionais do Play4Change, organizados por área funcional.

\begin{longtable}{p{1.3cm}p{3.5cm}p{8cm}}
\caption{Requisitos funcionais do sistema Play4Change} \label{tab:rf} \\
\hline
\textbf{ID} & \textbf{Área} & \textbf{Descrição} \\
\hline
\endfirsthead
\multicolumn{3}{l}{\small \textit{(continuação da Tabela~\ref{tab:rf})}} \\
\hline
\textbf{ID} & \textbf{Área} & \textbf{Descrição} \\
\hline
\endhead
\hline \multicolumn{3}{r}{\small \textit{continua na página seguinte}} \\
\endfoot
\hline
\endlastfoot
RF01 & Autenticação      & O sistema deve suportar autenticação para ambos os perfis de utilizador. \\
RF02 & Ingestão          & O Administrador deve poder submeter documentos PDF como fonte de conhecimento. \\
RF03 & Ingestão          & O Administrador deve poder submeter URLs como fonte de conhecimento. \\
RF04 & IA                & O sistema deve gerar automaticamente tópicos a partir do conteúdo submetido pelo Administrador. \\
RF05 & IA                & O sistema deve gerar automaticamente desafios (questões) a partir dos tópicos identificados. \\
RF06 & RAG               & O sistema deve armazenar o conhecimento extraído em forma de \textit{embeddings} e recuperá-lo para contextualizar a geração de conteúdo pela IA. \\
RF07 & RAG               & A regeneração completa de um tópico pelo Administrador deve invalidar e recalcular os \textit{embeddings} associados; a edição pontual de uma tarefa individual pode aceitar um vetor desactualizado (ver Secção~\ref{sec:invalidacao}). \\
RF08 & Desafios          & O sistema deve atribuir um desafio diário a cada Learner, de acordo com o seu percurso de progressão. \\
RF09 & Desafios          & O Learner deve receber \textit{feedback} imediato após responder a cada desafio. \\
RF10 & Consolidação      & O sistema deve implementar um percurso de remediação (\textit{struggle path}) que gere sub-tarefas adaptadas ao padrão de erro do Learner, reutilizando contextos de dificuldade semelhantes já resolvidos. \\
RF11 & Cliente Web       & O sistema deve disponibilizar um painel de administração acessível via browser. \\
RF12 & Cliente Web       & O sistema deve disponibilizar uma \textit{landing page} pública com informação sobre a plataforma e ligação para descarregar a aplicação móvel. \\
RF13 & Cliente Móvel     & O sistema deve disponibilizar uma aplicação móvel para os Learners, centrada na experiência de um desafio diário. \\
\end{longtable}

% -----------------------------------------------------------
\section{Requisitos Não Funcionais} \label{sec:rnf}
% -----------------------------------------------------------

Os requisitos não funcionais definem as propriedades e restrições do sistema que não dizem respeito a comportamentos específicos, mas à qualidade global da solução~\cite{bass2021}.
A Tabela~\ref{tab:rnf} apresenta os requisitos não funcionais identificados para o Play4Change.

\begin{table}[H]
\centering
\caption{Requisitos não funcionais do sistema Play4Change}
\label{tab:rnf}
\begin{tabular}{p{1.3cm}p{3cm}p{8.5cm}}
\hline
\textbf{ID} & \textbf{Categoria} & \textbf{Descrição} \\
\hline
RNF01 & Segurança       & O sistema deve proteger os utilizadores contra ataques \textit{Cross-Site Scripting} (XSS), removendo ou codificando HTML potencialmente malicioso nos conteúdos gerados por IA antes de os persistir ou renderizar. \\
RNF02 & Segurança       & O sistema deve aplicar \textit{rate limiting} por endereço IP para prevenir abusos nos endpoints públicos. \\
RNF03 & Segurança       & O sistema deve aplicar \textit{rate limiting} por endereço IP nas chamadas à API de IA, para controlo de custos e prevenção de abusos. \\
RNF04 & Segurança       & A comunicação entre clientes e servidor deve ser efetuada exclusivamente via HTTPS. \\
RNF05 & Usabilidade     & A aplicação móvel deve seguir os princípios de \textit{mobile-first design}, com interações mínimas para completar o desafio diário. \\
RNF06 & Escalabilidade  & A arquitetura do sistema deve permitir o aumento do número de utilizadores sem alterações estruturais ao domínio ou à aplicação; componentes de estado local por processo (e.g.\ o \textit{cache} de \textit{rate limiting}) são uma limitação conhecida a endereçar antes de escalar horizontalmente. \\
RNF07 & Manutenibilidade & O código deve seguir os princípios SOLID e ser coberto por testes unitários e de integração. \\
\hline
\end{tabular}
\end{table}

% -----------------------------------------------------------
\section{Casos de Uso Principais} \label{sec:casosdeuso}
% -----------------------------------------------------------

Os casos de uso descrevem as interações entre os atores e o sistema para atingir um objetivo específico~\cite{jacobson1992}.
São apresentados de seguida os três casos de uso centrais do Play4Change.

\subsection{CU01 --- Submeter Conteúdo e Gerar Desafios}

\textbf{Ator principal:} Administrador.\\
\textbf{Pré-condição:} O Administrador está autenticado no painel web.\\
\textbf{Fluxo principal:}
\begin{enumerate}
    \item O Administrador define o tópico e submete um documento PDF ou um URL como fonte de conhecimento; o sistema cria de imediato o tópico em estado \texttt{PENDING}.
    \item De forma assíncrona, o sistema extrai e limpa o texto do documento ou página.
    \item O sistema divide o texto em \textit{chunks} e invoca o modelo de IA para gerar módulos e os respetivos desafios (\textit{task templates}).
    \item Por cada desafio gerado, o sistema calcula o \textit{embedding}, verifica duplicação semântica face aos desafios já existentes no módulo, e persiste o resultado.
    \item Concluída a geração, o tópico transita automaticamente para o estado \texttt{ACTIVE}, sem aprovação explícita do Administrador; este pode editar um desafio específico ou desencadear a regeneração completa do tópico.
\end{enumerate}
\textbf{Pós-condição:} O tópico está \texttt{ACTIVE}, com os desafios e o conhecimento associado indexados e disponíveis para atribuição aos Learners.

\subsection{CU02 --- Responder ao Desafio Diário}

\textbf{Ator principal:} Learner.\\
\textbf{Pré-condição:} O Learner está autenticado na aplicação móvel e tem um desafio diário atribuído.\\
\textbf{Fluxo principal:}
\begin{enumerate}
    \item O Learner abre a aplicação móvel e visualiza o desafio do dia.
    \item O sistema apresenta o desafio correspondente ao seu percurso atual, já gerado antecipadamente pela IA durante o pipeline de ingestão do tópico --- não há recuperação vetorial em tempo real neste passo.
    \item O Learner submete a sua resposta.
    \item O sistema avalia a resposta e apresenta \textit{feedback} imediato, com explicação contextualizada.
    \item O sistema regista o resultado e atualiza o perfil de progressão do Learner.
\end{enumerate}
\textbf{Pós-condição:} O resultado é registado; se a resposta estiver incorreta, é desencadeada a sessão de remediação descrita em CU03.

\subsection{CU03 --- Remediação Adaptativa (Struggle Path)}

\textbf{Ator principal:} Learner.\\
\textbf{Pré-condição:} O Learner respondeu incorretamente a um desafio do dia.\\
\textbf{Fluxo principal:}
\begin{enumerate}
    \item O sistema classifica deterministicamente o padrão do erro (por exemplo, atraso na submissão ou opção adjacente à correta).
    \item De forma assíncrona, o sistema gera três sub-tarefas mais simples, centradas no padrão de erro detetado, reutilizando um ramo já gerado para um contexto semelhante sempre que a similaridade semântica o permita.
    \item O Learner resolve as sub-tarefas; se as completar todas corretamente, a sessão é marcada como resolvida e o desafio original do dia é reposto para nova tentativa.
    \item Se o Learner voltar a falhar após esgotar a profundidade de remediação definida, o sistema substitui as sub-tarefas por uma sessão de explicação sintetizada por IA, com possibilidade de diálogo de esclarecimento.
\end{enumerate}
\textbf{Pós-condição:} O Learner é conduzido de volta ao desafio original, com uma remediação centrada no seu erro específico; esta adaptação ocorre dentro do próprio desafio do dia, não altera a escolha de que tópico apresentar em dias seguintes --- essa seleção continua a ser feita pelo próprio Learner, através de inscrição explícita nos tópicos disponíveis.
